\section{A Second Round of Improvements}

Ever since we started CDCL, I have been promising massive speedups over DPLL 
and yet our SAT solver has gotten slower each iteration. Before you treat me 
like Jesus and crucify me,
\footnote{My friends told me that it was too funny not to include so it stays.}
I knew and deep down you knew that the way we were implementing was not optimal
to begin with. So in a small way, this is your fault too. 

\subsection{Removing Hashing}

The first thing we're going to do is remove the maps and instead use arrays 
instead so we don't have to deal with hashing. Our interface becomes like so:

\begin{lstlisting}
    public interface:
        SATSolver(string input)
        bool solveCNF()

    private interface:
        ...
        int[][] var_to_clause
        (int, int)[] clause_to_var
\end{lstlisting}

In case you need a reminder as it has been a while, these are the data structures 
used to make watched literals happen. Before they were maps where the keys were variables 
or clause indices, but now, each index in the arrays are the keys and the values at those 
indices are, well, the values. There is a small implementation detail concerning 
\textit{clause\_to\_var}, which is that now that the number of clauses are not fixed, 
everytime we add in a new clause, we need to add in a new entry to \textit{clause\_to\_var}
before we let our \textit{createWatched(int, int)} overload handle assigning the values 
to that entry. Basically, to avoid going out of bounds for the data structure, we need to 
change the bounds before we access the new boundary. While we could let 
\textit{createWatched(int, int)} handle the allocation as well, it would make it less 
versatile as it would have to only operate on the last clause in the database instead of 
any clause in the database.

The second point to mention is that now our propagation logic doesn't work because we 
need to create a mapping from literals (positive and negative) to indices (nonnegative).
We use the following mapping:
$$ 2 * abs(lit) + (lit < 0) $$
This is how the mapping works out (left column is literal value, right column is index):
\begin{FVerbatim}
 1 → 0
-1 → 1
 2 → 2
-2 → 3
 3 → 4
-3 → 5
  ...
\end{FVerbatim}

The last step is to remove the duplicates set we use in clause resolution.
We're going to add in an array of booleans to our SAT solver where index $i$ 
is whether or not we've seen that variable during our clause resolution.
However, since our (my) implementation of first UIP clause resolution is 
such a big fixer-upper, it deserves its own subsection.

\subsection{Fixing Clause Resolution}

Our current implementation of clause resolution is pretty long and inefficient. This was 
written before we understood implication graphs fully. Now we do understand implication graphs, 
so let's exploit that structure.

The idea is this: instead of rescanning over and over to make sure our resolution clause
has just one current level variable, we'll keep track of which variables we've seen as 
well as which ones are from this decision level. Then, as we perform clause resolution, 
we'll update the seen and count variables to reflect the current state of the resolution 
clause, then once we've encountered our stopping condition, we'll use which variables 
we've kept marked as seen to construct our learned clause directly. This is very different 
from the way we've been performing clause resolution, which involved repeatedly scanning the 
resolution clause and recomputing how far along we are in the UIP process. 

How are we using the implication graph in our new method? Well since the implication 
graph is an acyclic graph, that means we can use its structure for our UIP traversal. 
There is a relationship between the implication graph and our trail, in that our trail 
is just a compressed version of the implication graph, meaning it doesn't keep track of 
which variable assignments led to which propagations. However, they both follow the same 
structure in that the edge of our trail are the edges of our implication graphs and the 
order of variables in the trail and graphs are the same. Because the implication graph 
is acyclic, once we encounter a variable during clause resolution, we will never encounter 
it again while we're resolving. Thanks to that property, we can traverse the trail backwards 
without having to worry about going over the same parts of the trail/graph, which was 
something that we were sort of doing in the old method when we rescanned the whole 
conflict clause, whereas now, we are only concerned with the differences to the resolution
clause made across each iteration. 

Our resolution algorithm becomes as follows:
\begin{algorithm}[H]
\caption{Improved First UIP}\label{alg:34}
\begin{algorithmic}[1]
\Procedure{resolveConflict}{conflict[1..n]}

\State levelCount $\gets$ 0
\State curLevel $\gets$ decisionLevel
\For{each $var$ in $conflict$}
    \If{seen[\Call{abs}{var}] = false}
        \State seen[\Call{abs}{var}] $\gets$ true 
        \State tIdx $\gets$ varToTrail[\Call{abs}{var}]
        \If{trail[tIdx].decisionLevel = curLevel}
            \State levelCount $\gets$ levelCount + 1
        \EndIf
    \EndIf
\EndFor

\For{each $entry$ in $trail$ from last} \Comment traverse the trail in reverse order
    \State idx $\gets$ entry.lit

    \If{seen[idx] = true}
        \State reason $\gets$ Clauses[entry.reasonIdx]

        \For{each $var$ in $reason$}
            \If{seen[\Call{abs}{var}] = false}
                \State seen[\Call{abs}{var}] $\gets$ true
                \State tIdx $\gets$ varToTrail[\Call{abs}{var}]
                \If{trail[tIdx].decisionLevel = curLevel}
                    \State levelCount $\gets$ levelCount + 1
                \EndIf
            \EndIf
        \EndFor

        \State seen[idx] $\gets$ false
        \State levelCount $\gets$ levelCount - 1
    \EndIf

    \If{levelCount = 1}
        \State break
    \EndIf
\EndFor

\State learned[1..m]
\For{i $\gets$ 1 up to numVars}
    \If{seen[i] = true}
        \State lit $\gets$ trail[varToTrail[i]].lit
        \State $\Call{append}{learned, -lit}$
        \State seen[i] $\gets$ false
    \EndIf
\EndFor

\State return learned

\EndProcedure
\end{algorithmic}
\end{algorithm}

\noindent
Here's the rundown of what's going on. In the first block, we're marking which 
variables we've seen and how many of them are from the current decision level from 
the initial conflicting clause. Next, we start traversing our trail from the most recent 
trail entry, and we move towards the front of the trail. If we've seen the variable from 
the current trail entry, we grab its reason clause, and for each new unseen variable, 
we mark it as seen and increment the current decision level count. Next, we ``remove''
the variable we just processed by unmarking it as seen and decrementing \textit{levelCount}.

Once we have only one literal from the current level, we construct a new learned clause 
by iterating over the entire \textit{seen} array and we add the negative literal to the 
learned clause because the learned clause is supposed to learn the opposite of what we've 
done so we don't repeat it. Then we unmark everything we have seen so that way we can reuse 
the same array the next time there's a conflict instead of allocating a new one each time. 

The results for this are very good, but I won't tell you what they are, because there's 
one more thing I want to do. You see, we iterate over the entire \textit{seen} array,
which is as large as the number of variables in our equation. However, it's very unlikely 
that we'll see many variables because first UIPs shrink how large the learned clause is, 
so there's less to be seen in the first place, but also because the more we learn, generally,
the learned clause gets smaller because each new learned clause adds more constrictions on 
which assignments are valid, and as we constrict which ones are valid and which ones aren't,
we inch towards learning the empty clause signaling that nothing is valid. This is just a 
long winded way to say that the iterating over the entire seen array each time, while O(n),
could be better.

But isn't that microoptimizing?
\footnote{Microoptimizing is like cutting your hair to help you reach a weight loss goal, 
instead of fixing your running form so you'll be able to have better cardio.}
After all, our two test sets have 20 and 50 variables 
respectively, which isn't much in the grand scheme of things. And you're right, our test 
sets are pretty small. However, the goal is to build a SAT solver that would scale well 
and most SAT solvers can handle problems with millions of variables pretty well, so I'd 
like this one to too. Making it scale well isn't the same as microoptimizing. Microoptimizing 
would be like saying \textit{++i} over \textit{i++}.
\footnote{Pre vs post incrementation in loops generally does not affect runtime and whether 
or not it actually matters is heavily debated.}

In order to do this, we're going to keep track of the indices of which variables we've
seen in an array and then iterate over that array. We will also need to keep track of 
which indices are in which position so we can swap and pop for efficient removal of 
elements.
\footnote{I love swap and pop if only for the name. I encourage everyone to do swap and 
pops so we can say ``swap and pop'' more than we currently say ``swap and pop.''}
Our object now looks like this:

\begin{lstlisting}
    public interface:
        SATSolver(string input)
        bool solveCNF()

    private interface:
        ...
        int[][] var_to_clause
        (int, int)[] clause_to_var

        bool[] seen
        int[] seenIdx
        int[] seenPos
\end{lstlisting}

\noindent
Here a subtle detail about these data structures. \textit{seen} and \textit{seenPos}
are a fixed size of \textit{numVars} because they require a full mapping from variable
to value at all times, but \textit{seenIdx} is dynamically sized because the number of 
variables that we see each clause resolution changes. So if you are implementing this 
yourself, make sure that \textit{seen} and \textit{seenPos} are allocated enough space 
upon initialization. Personally, I also reserve \textit{numVars} amount of space in 
\textit{seenIdx} so that it never has to reallocate space while solving, but it's not 
strictly necessary. 

This is how these data structures change the algorithm:

\begin{algorithm}[H]
\caption{Improved-improved First UIP}\label{alg:35}
\begin{algorithmic}[1]
\Procedure{resolveConflict}{conflict[1..n]}

\State levelCount $\gets$ 0
\State curLevel $\gets$ decisionLevel
\For{each $var$ in $conflict$}
    \If{seen[\Call{abs}{var}] = false}
        \State vIdx $\gets$ \Call{abs}{var}
        \State seen[vIdx] $\gets$ true 

        \State \Call{append}{seenIdx, vIdx}
        \State seenPos[vIdx] $\gets$ \Call{size}{seenIdx}

        \State tIdx $\gets$ varToTrail[vIdx]
        \If{trail[tIdx].decisionLevel = curLevel}
            \State levelCount $\gets$ levelCount + 1
        \EndIf
    \EndIf
\EndFor

\For{each $entry$ in $trail$ from last} \Comment traverse the trail in reverse order
    \State idx $\gets$ entry.lit
    
    \If{seen[idx] = true}
        \State reason $\gets$ Clauses[entry.reasonIdx]
        \For{each $var$ in $reason$}
            \State vIdx $\gets$ \Call{abs}{var}

            \If{seen[\Call{abs}{var}] = false}
                \State seen[vIdx] $\gets$ true 

                \State \Call{append}{seenIdx, vIdx}
                \State seenPos[vIdx] $\gets \Call{size}{seenIdx}$

                \State tIdx $\gets$ varToTrail[vIdx]
                \If{trail[tIdx].decisionLevel = curLevel}
                    \State levelCount $\gets$ levelCount + 1
                \EndIf
            \EndIf
        \EndFor

        \State seen[idx] $\gets$ false

        \State pos $\gets$ seenPos[idx]
        \State last $\gets$ seenIdx.last
        \State seenIdx[pos] $\gets$ last
        \State seenIdx[last] $\gets$ pos

        \State \Call{removeLast}{seenIdx}

        \State levelCount $\gets$ levelCount - 1
    \EndIf

    \If{levelCount = 1}
        \State break
    \EndIf
\EndFor

\State learned[1..m]
\While{seenIdx is not empty}
    \State lIdx $\gets$ seenIdx.last
    \State lit $\gets$ trail[varToTrail[lIdx]].lit
    \State $\Call{append}{learned, -lit}$
    \State seen[lIdx] $\gets$ false
    \State \Call{removeLast}{seenIdx}
\EndWhile

\State return learned

\EndProcedure
\end{algorithmic}
\end{algorithm}

In this code, not much changes other than that we index using the seenIdx and seenPos
data structures to achieve a quicker deletion and modification of the learned clause
by both swapping and popping seenIdx as well as seenPos.

\subsection{Improvements - Results}

I'll let the results speak for themselves:
\noindent
uf20.91: \textbf{264ms}

\noindent
UUF50.218.1000: \textbf{2385ms}

\begin{FVerbatim}
On all 1000 equations in uf20-91:

Improved CDCL:             284ms (~0.28 seconds).
First UIPs:                406ms (~0.41 seconds).
Backjumping:               380ms (~0.38 seconds).
Clause Learning:           425ms (~0.43 seconds).
Iterative DPLL (AbP):      328ms (~0.33 seconds).
Iterative DPLL (PbA):      324ms (~0.32 seconds).
Watched Literals:          349ms (~0.35 seconds).
Pure Literal Elimination:  1548ms (~1.55 seconds).
Unit Propagation:          868ms (~0.87 seconds).
Recursive Backtracking:    1,540,613ms (~25.68 minutes).
Brute Force:               1,770,851ms (~29.51 minutes).
Short-Circuited (Basic):   49,874ms (~50 seconds).
Parallel:                  10,589ms (~10.5 seconds).

On all 1000 equations in UUF50.218.1000:

Improved CDCL:            2,385ms (~2.39 seconds).
First UIPs:               14,012ms (~14.01 seconds).
Backjumping:              11,291ms (~11.29 seconds).
Clause Learning:          9,652ms (~9.65 seconds).
Iterative DPLL (AbP):     6,369ms (~6.37 seconds).
Iterative DPLL (PbA):     6,365ms (~6.37 seconds).
Watched Literals:         7,130ms (~7.13 seconds).
Pure Literal Elimination: 231,107ms (~3.85 minutes).
Unit Propagation:         140,041ms (~2.34 minutes).
\end{FVerbatim}
